\documentclass[../numerical_methods.tex]{subfiles}
\begin{document}
\section{Aula 18 - 11 de Junho, 2026}
\subsection{Motivações}
\begin{itemize}
	\item Critério de Parada do Método da Potência;
	\item Potências Inversas;
\end{itemize}
\subsection{Método da Potência com Critério de Parada e Potência Inversa}
Dado \(\varepsilon > 0\) e \(k_{max}\in \mathbb{N}\) um valor máximo de iterações, queremos que o método das potências pare quando \(k=k_{max}\), ou se
\[
	| \lambda_{1}^{(k+1)} - \lambda_{1}^{(k)} | < \varepsilon,
\]
ou mesmo utilizando o teste de alinhamento entre os vetores \(\underline{y}^{(k+1)}\) e \(\underline{y}^{(k)}\), que seria quando o cosseno do ângulo entre eles fica aproximadamente 1:
\[
	| | \underline{y}^{(k+1)} \cdot \underline{y}^{(k)} | - 1 | < \varepsilon .
\]
Usando o critério de parada, podemos fazer uma implementação inicial do método das potências:
\begin{listing}[H]
	\caption{Método da Potência com critério de parada}
	\begin{minted}[bgcolor = DarkBackground, linenos=true, breaklines]{Python}
import numpy as np
from scipy import linalg as lin

# Método da Potência (Maior Autovalor em Módulo)
def potencia(A, tol):
    k = 0
    kmax = 10000
    erro = np.inf
    n = A.shape[0]

    # Vetor inicial chuta 1 na primeira posição
    y0 = np.zeros(n)
    y0[0] = 1.0

    while erro > tol and k < kmax:
        x = A.dot(y0)
        y = x / np.linalg.norm(x)
        erro = np.abs(np.abs(y0.dot(y)) - 1)
        y0 = y
        k += 1

    lambda_ = y.dot(A.dot(y))
    return lambda_, y, k
   \end{minted}
\end{listing}

Agora, seja \(A\in \mathbb{M}_{n}\) uma matriz invertível e suponha que os seus autovalores estão ordenados da seguinte forma:
\[
	| \lambda_1 | \geq | \lambda_2 | \geq \dotsc > \lambda_{n} > 0.
\]
Com isso, se calcularmos o menor autovalor em módulo \(\lambda_{n}\) e o seu autovetor associado \(\underline{v}_{n}\), obteremos
\[
	A \underline{v}_{j} = \lambda_{j}\underline{v}_{j} \Longleftrightarrow A^{-1}\underline{v}_{j} = \frac{1}{\lambda_{j}}\underline{v}_{j},
\]
mostrando que os autovalores de \(A^{-1}\) são dados pelas recíprocas dos autovalores de A; logo, \(\frac{1}{\lambda_{n}}\) é o autovalor dominante de \(A^{-1}\)! Isso permite repetirmos os raciocínios do método da potência normal da seguinte forma: dado um chute inicial \(\underline{x}^{(0)}\in \mathbb{R}^{n}\) não nulo, comece computando \(\underline{y}^{(0)}=\frac{\underline{x}^{(0)}}{\Vert \underline{x}^{(0)} \Vert_{2}}\). A partir disso, para os k subsequentes, faça:
\[
	\boxed{\underline{x}^{(k)} = A^{-1}\underline{y}^{(k-1)},\; \underline{y}^{(k)} = \frac{\underline{x}^{(k)}}{\Vert \underline{x}^{(k)} \Vert_{2}},\; \lambda_{n}^{(k)} = \underline{y}^{(k)}\cdot A \underline{y}^{(k)}}
\]
Apesar de ser bastante parecido com o anterior, é necessário resolver o sistema linear \(A\underline{x}^{(k)} = \underline{y}^{(k-1)}\) a cada iteração, que pode ser feito calculando a decomposição \(LU\) de A e armazenando \(L\) e \(U\), seguido dos algoritmos de substituição para resolver os sistemas obtidos. No entanto, como calcular um autovalor arbitrário \(\lambda_{j}\) e o seu autovetor associado \(\underline{v}_{j}\)?

A ideia é considerar \(\alpha \in \mathbb{C}\) que não seja um autovalor de \(A\), de modo que
\[
	\Lambda (A) = \{\lambda_1, \dotsc , \lambda_{n}\} \Rightarrow \Lambda (A-\alpha I)=\{\lambda_{1}-\alpha , \dotsc , \lambda_{n}-\alpha \};
\]
assim, os autovalores de \(B = (A-\alpha I)^{-1}\) são da forma
\[
	\boxed{\gamma_{j}=\frac{1}{\lambda_{j}-\alpha }}
\]
e, portanto, quanto mais próximo \(\alpha \) estiver de \(\lambda_{j}\), mais dominante será \(\gamma_{j}\)! Por isso, o valor \(\alpha \) leva o nome de \textbf{deslocamento}. Então, como resposta à pergunta prévia, basta aplicar o método da potência inversa na matriz deslocada \(B\).
A implementação no código pode ser feita com
\begin{listing}[H]
	\caption{Implementação do Método da PotÊncia Inversa}
	\begin{minted}[bgcolor = DarkBackground, linenos=true, breaklines]{Python}
# Método da Potência Inversa (Menor Autovalor em Módulo)
def potencia_inv(A, tol):
    k = 0
    kmax = 10000
    erro = np.inf
    n = A.shape[0]
    y0 = np.zeros(n)
    y0[0] = 1.0

    # Fatoração LU para evitar resolver o sistema do zero em todo loop
    lu, piv = lin.lu_factor(A)

    while erro > tol and k < kmax:
        # Resolve usando a matriz já fatorada (muito mais rápido)
        x = lin.lu_solve((lu, piv), y0)
        y = x / np.linalg.norm(x)
        erro = np.abs(np.abs(y0.dot(y)) - 1)
        y0 = y
        k += 1

    lambda_ = y.dot(A.dot(y))
    return lambda_, y, k
\end{minted}
\end{listing}
Para testar, usamos
\begin{listing}[H]
	\begin{minted}[bgcolor = DarkBackground, linenos=true, breaklines]{Python}
# --- Execução do Exemplo e Validação ---

A = np.array([[-4, 14, 0],
              [-5, 13, 0],
              [-1, 0, 2]], dtype='double')

# 1. Cálculo nativo do Python (Gabarito)
D_python, V_python = np.linalg.eig(A)

print("=" * 50)
print("AUTOVALORES REAIS (Nativo do Python):")
print(D_python)
print("=" * 50)

# 2. Teste do Método da Potência Regular
print('\nCálculo do MAIOR autovalor em módulo de A')
lambda_, y, k = potencia(A, 0.000001)
print('-> Método da Potência: %.4f (usando %d iterações)' % (lambda_, k))
print('-> Gabarito do Python: %.4f' % (np.max(np.abs(D_python))))

# 3. Teste do Método da Potência Inversa
print('\nCálculo do MENOR autovalor em módulo de A')
lambda_inv, y_inv, k_inv = potencia_inv(A, 0.000001)
print('-> Potência Inversa:  %.4f (usando %d iterações)' % (lambda_inv, k_inv))
print('-> Gabarito do Python: %.4f' % (np.min(np.abs(D_python))))
print("=" * 50)
\end{minted}
\end{listing}

\begin{listing}[H]
	\begin{minted}[bgcolor = DarkBackground, linenos=true, breaklines]{Python}
# ==========================================
# DEFINIÇÃO DE MATRIZ DE EXEMPLO (10x10)
# ==========================================

B = np.array([
    [  1,   0,   0,   0,   0,   0,   0,   0,   0,   0],
    [ -3,   2,   0,   0,   0,   0,   0,   0,   0,   0],
    [  3,  -1,   3,   0,   0,   0,   0,   0,   0,   0],
    [ -6,   3,  -3,   4,   0,   0,   0,   0,   0,   0],
    [ -4,  -1,   1,   1,   5,   0,   0,   0,   0,   0],
    [  0,   6,   6,  -3,  -3,   6,   0,   0,   0,   0],
    [-44,  25,  -7,   2,  -7,   1,   7,   0,   0,   0],
    [-45,   7, -19,  14,   6,  -4,   0,   8,   0,   0],
    [153, -20,  72, -32, -31,  11,   4,   1,   9,   0],
    [  1,   5,  -5,   0, -15,  -3,   8,   3,   1,  10]
], dtype='double')

# Cálculo dos autovalores de referência (gabarito do Python)
D_python_B, _ = np.linalg.eig(B)

# ==========================================
# TESTES COM A MATRIZ B
# ==========================================

print("\n" + "=" * 60)
print("TESTES COM A MATRIZ B")
print("=" * 60)
print("Todos os autovalores reais de B (Nativo):\n", np.round(D_python_B, 4))
print("-" * 60)

# Maior de B
print('Cálculo do MAIOR autovalor em módulo de B')
lambda_B_max, _, k = potencia(B, 0.000001)
print('-> Potência: %.4f usando %d iterações' % (lambda_B_max, k))
print('-> Python:   %.4f\n' % (np.max(np.abs(D_python_B))))

# Menor de B
print('Cálculo do MENOR autovalor em módulo de B')
lambda_B_min, _, k = potencia_inv(B, 0.000001)
print('-> Potência Inversa: %.4f usando %d iterações' % (lambda_B_min, k))
print('-> Python:           %.4f' % (np.min(np.abs(D_python_B))))
print("=" * 60)
\end{minted}
\end{listing}

O problema é que a resposta levantou outra pergunta: como escolher \(\alpha \)? Como \(\alpha \) mora no plano dos complexos, usaremos um pouco deles para conseguir.
\begin{def*}
	Seja \(A\in \mathbb{M}_{n}\); o \textbf{disco de Gershgorin} de centro \(a_{ii}\) e raio \(r_{i}\) associado à i-ésima linha de A é definido como
	\[
		\mathcal{D}_{A}(a_{ii}, r_{i}) = \{z\in \mathbb{C}:\; | z-a_{ii} | \leq r_{i}\},  \quad r_{i} = \sum\limits_{\substack{j=1 \\ j\neq i}}^{n} | a_{ij} |
	\]
\end{def*}
Com ele, podemos propor uma proposição que fornece um teste de rejeição no cálculo dos autovalores e uma boa região, apesar de grande, para escolher o deslocamento \(\alpha \) no método da potência inversa:
\begin{prop*}
	Seja \(A\in \mathbb{M}_{n}\); então,
	\[
		\Lambda (A) \subseteq \biggl(\bigcup_{i=1}^{n}\mathcal{D}_{A}(a_{ii}, r_{i})\biggr)\cap \biggl(\bigcup_{i=1}^{n}\mathcal{D}_{A^{T}}(a_{ii}, r_{i})\biggr).
	\]
\end{prop*}
Nos casos onde \(A=A^{T}\), o teorema espectral garante que \(\Lambda (A)\subseteq \mathbb{R}\), permitindo restringir de um disco a um intervalo!
\subsection{Alguns Conceitos Adicionais}
Para ajudar no Trabalho 3, sobre Ordenação de Páginas, alguns conceitos adicionais são necessários.
\begin{def*}
	Um vetor \(\underline{p}\in \mathbb{R}^{n}\) com \(p_{i}\geq 0\) e \(\sum\limits_{i=1}^{n}p_{i}=1\) é chamado \textbf{vetor estocástico}.

	Uma matriz cujos vetores coluna são estocásticos é chamada \textbf{matriz estocástica}. \(\square\)
\end{def*}
\begin{prop*}
	Se \(A\in \mathbb{M}_{n}\) e \(\underline{p}\in \mathbb{R}^{n}\) são estocásticos, então o vetor \(A \underline{p}\) é estocástico e \(\lambda =1\) é um autovalor de \(A\).
\end{prop*}
Dada uma matriz estocástica A e uma distribuição de probabilidade \(\underline{p}^{(t)}\) no tempo t, podemos descobrir a distribuição no tempo (t+k) pelo processo
\[
	\underline{p}^{(t+k)}=A^{k}\underline{p}^{(t)}.
\]
\begin{example}
	Se uma cidade tem pessoas que ficam apenas nos estados corporais magra (1), gorda (2) e sarada (3), com distribuição inicial \(\underline{p}^{(0)} = (1/3, 1/3, 1/3)^{T}\), a matriz estocástica determina como esse sistema evoluirá com o passar dos tempos; por exemplo:
	\[
		A = \begin{bmatrix}
			\frac{1}{2} & 0           & \frac{1}{4} \\
			0           & \frac{1}{2} & \frac{1}{4} \\
			\frac{1}{2} & \frac{1}{2} & \frac{1}{2}
		\end{bmatrix} \Rightarrow \underline{p}^{(1)} = A \underline{p}^{(0)}  = \begin{bmatrix}
			\frac{1}{4} \\
			\frac{1}{4} \\
			\frac{1}{2}
		\end{bmatrix}.
	\]
\end{example}
\begin{theorem*}[Perron-Frobenius]
	Seja \(A\in \mathbb{M}_{n}\) uma matriz estocástica; então,
	\begin{itemize}
		\item[1)] 1 é o autovalor dominante de A; e
		\item[2)] O autovetor \(\underline{v}\) associado a \(\lambda \) possui todas as entradas ou positivas, ou negativas; em particular, para \(\lambda \), existe um único autovalor que é estocástico.
	\end{itemize}
\end{theorem*}
Portanto, a convergência do processo de Markové garantida pelo método das potências! O autovetor \(\underline{v}\) tal que \(\underline{p}^{(k)}\to \underline{v}\) é chamado \textbf{vetor estacionário} de A.

\end{document}
